\section{Future Work}
\label{sec:futurework}

The present results characterise a first working point of the hybrid design
(1.45B total / $\sim$768M active parameters, $\sim$6.3B training tokens). The
following steps are needed before any claim about ``useful agent'' capability, or
about the architectural choices themselves, can be substantiated.

\paragraph{Scaling pretraining compute.}
Training consumed roughly $6.3$B tokens (one epoch over $6.1$M $\times$ 1024-token
chunks), which is well below compute-optimal for a model of this size. Scaling to
$10^2$--$10^3$B tokens is expected to be the single largest source of improvement and
is the natural next run.

\paragraph{Matched-baseline comparison.}
The current evaluation does not isolate the contribution of the hybrid architecture.
A controlled study at an identical token budget should compare against (i) a dense
Transformer of matched active parameters and (ii) a public same-scale base model
(e.g.\ Pythia-1.4B) at matched tokens, before attributing any gain to MLA, Mamba2, or
the aux-loss-free MoE.

\paragraph{Post-training.}
So far only next-token-prediction pretraining and a preliminary SFT stage have been
run. Supervised instruction tuning followed by preference optimisation (DPO/RLHF)
should be evaluated end-to-end, since the current gap is dominated by missing
preference learning rather than by the architecture.

\paragraph{Evaluation coverage.}
Beyond C-Eval, TruthfulQA, and GSM8K, the agent setting requires code generation
(HumanEval/MBPP), long-context retrieval, and multi-turn tool-use tasks. Every reported
score should be accompanied by the token budget and by matched-baseline deltas.

\paragraph{Architecture ablations.}
The routed-expert count and top-$k$, the MLA KV-compression rank, and the
Mamba2/attention interleaving ratio should be swept under a fixed token budget, so that
each component's contribution is measured rather than assumed.

\paragraph{Long context.}
The current context window is 1024 tokens. Extending it via position interpolation and
continued pretraining is a prerequisite for the document- and tool-heavy workloads the
architecture targets.

\paragraph{Inference efficiency.}
The constant-size KV cache and SSM state, together with the TileLang kernel path, should
be benchmarked on throughput and peak memory, in order to quantify the active-parameter
efficiency claim on realistic agent workloads.
